% Page de garde : tout est placé en coordonnées absolues (cm) depuis le coin bas gauche.
\begin{titlepage}
\thispagestyle{empty}
\begin{tikzpicture}[remember picture,overlay,x=1cm,y=1cm]
\begin{scope}[shift={(current page.south west)}]

% --- bloc bleu supérieur, coupé en diagonale ---
\shade[top color=bleuENP!70!black, bottom color=bleuENP]
  (0,29.7) -- (21,29.7) -- (21,15.2) -- (0,12.6) -- cycle;
\fill[orangeTP] (0,12.6) -- (21,15.2) -- (21,14.95) -- (0,12.35) -- cycle;

% --- champ d'ondes progressives en filigrane ---
\begin{scope}
  \clip (0,29.7) -- (21,29.7) -- (21,15.2) -- (0,12.6) -- cycle;
  \foreach \k in {0,...,13} {
    \draw[white, opacity=0.05+0.012*\k, line width=0.5pt, smooth, domain=0:21, samples=120]
      plot (\x, {14.2+1.05*\k + (0.25+0.03*\k)*sin(\x*(38+4*\k) + 37*\k)});
  }
\end{scope}

% --- grand « 01 » en transparence ---
\node[anchor=north east, white, opacity=0.10, inner sep=0, font=\sffamily\bfseries\fontsize{170}{170}\selectfont]
  at (20.4,29.2) {01};

% --- logo sur carte blanche ---
\fill[white, rounded corners=6pt] (1.5,24.6) rectangle (6.2,28.7);
\node at (3.85,26.65) {\includegraphics[width=4.0cm]{figures/logo_enp.png}};

% --- surtitre et titre ---
\node[anchor=west, orangeTP, font=\sffamily\bfseries\small] at (1.5,23.2)
  {\textls[200]{TRAVAUX PRATIQUES \,·\, RÉSEAUX ÉLECTRIQUES}};
\node[anchor=west, white, font=\sffamily\bfseries\fontsize{50}{50}\selectfont] at (1.35,21.5)
  {Propagation};
\node[anchor=west, white, font=\sffamily\fontsize{25}{28}\selectfont] at (1.5,19.75)
  {dans les lignes électriques};
\fill[orangeTP] (1.55,18.75) rectangle (4.6,18.9);

% --- signal : réponse en escalier de v2 à l'extrémité ouverte ---
\draw[orangeTP!40, line width=5pt, opacity=0.35, line join=round]
  (1.5,15.25) -- (4.0,15.25) -- (4.0,17.35) -- (6.2,17.35) -- (6.2,16.1) -- (8.4,16.1)
  -- (8.4,16.95) -- (10.6,16.95) -- (10.6,16.4) -- (12.8,16.4) -- (12.8,16.75) -- (15.0,16.75)
  -- (15.0,16.52) -- (17.2,16.52) -- (17.2,16.66) -- (19.5,16.66);
\draw[orangeTP, line width=1.6pt, line join=round]
  (1.5,15.25) -- (4.0,15.25) -- (4.0,17.35) -- (6.2,17.35) -- (6.2,16.1) -- (8.4,16.1)
  -- (8.4,16.95) -- (10.6,16.95) -- (10.6,16.4) -- (12.8,16.4) -- (12.8,16.75) -- (15.0,16.75)
  -- (15.0,16.52) -- (17.2,16.52) -- (17.2,16.66) -- (19.5,16.66);
\fill[white] (4.0,17.35) circle (2.2pt);
\node[anchor=south west, white, font=\sffamily\footnotesize, opacity=0.85] at (4.1,17.45)
  {$v_2 \approx (1+\Gamma)\,v_{\mathrm{inc}}$};
\node[anchor=north, white, font=\sffamily\scriptsize, opacity=0.8] at (5.1,16.0) {$2\tau$};
\draw[white, opacity=0.8, |<->|, line width=0.5pt] (4.0,16.3) -- (6.2,16.3);

% --- sous-titre ---
\node[anchor=west, text width=17.5cm, bleuENP!70!black, font=\large\itshape] at (1.5,10.2)
  {Modèle de Bergeron et règle des trapèzes : régimes transitoires\\ d'une ligne longue sans pertes};

% --- cartes d'information ---
\foreach \x/\titre/\ligneA/\ligneB in {
    1.5/RÉALISÉ PAR/Mohammed ZANE/Groupe 1,
    8.0/ENCADRÉ PAR/Prof. Karim SEBAA/Réseaux Électriques,
    14.5/ANNÉE/2026 -- 2027/TP n\textsuperscript{o}\,01} {
  \fill[bleuENP!6] (\x,5.5) rectangle ++(5.2,2.9);
  \fill[orangeTP] (\x,5.5) rectangle ++(0.14,2.9);
  \node[anchor=north west, bleuENP, font=\sffamily\bfseries\scriptsize] at (\x+0.45,8.05)
    {\textls[180]{\titre}};
  \node[anchor=north west, black, font=\sffamily\bfseries\large] at (\x+0.45,7.35) {\ligneA};
  \node[anchor=north west, black!65, font=\sffamily\small] at (\x+0.45,6.5) {\ligneB};
}

% --- pied de page ---
\draw[bleuENP!25, line width=0.4pt] (1.5,3.1) -- (19.5,3.1);
\node[anchor=west, black!60, font=\sffamily\small] at (1.5,2.5)
  {École Nationale Polytechnique \,·\, Département de Génie Électrique};
\node[anchor=east, black!60, font=\sffamily\small] at (19.5,2.5) {Alger};
\fill[bleuENP] (0,0) rectangle (21,0.55);
\fill[orangeTP] (0,0.55) rectangle (7,0.7);

\end{scope}
\end{tikzpicture}
\null
\end{titlepage}
